\documentclass[a4paper,14pt]{extarticle}

\usepackage[T2A]{fontenc}
\usepackage[utf8]{inputenc}
\usepackage[russian]{babel}
\usepackage{geometry}
\usepackage{setspace}
\usepackage{indentfirst}
\usepackage{array}
\usepackage{booktabs}
\usepackage{listings}
\usepackage{xcolor}

\geometry{
    left=30mm,
    right=15mm,
    top=20mm,
    bottom=20mm
}

\onehalfspacing
\setlength{\parindent}{1.25cm}
\setlength{\parskip}{0pt}
\setlength{\emergencystretch}{3em}

\definecolor{codebg}{RGB}{248,248,248}
\lstdefinestyle{cudastyle}{
    language=C++,
    basicstyle=\ttfamily\normalsize,
    keywordstyle=\color{blue!60!black},
    commentstyle=\color{green!40!black},
    stringstyle=\color{orange!50!black},
    backgroundcolor=\color{codebg},
    frame=single,
    breaklines=true,
    showstringspaces=false,
    tabsize=4
}

\begin{document}

\begin{titlepage}
    \begin{center}
        \textbf{МОСКОВСКИЙ АВИАЦИОННЫЙ ИНСТИТУТ}\\
        \textbf{(НАЦИОНАЛЬНЫЙ ИССЛЕДОВАТЕЛЬСКИЙ УНИВЕРСИТЕТ)}\\
        Институт №8 «Компьютерные науки и прикладная математика»
    \end{center}

    \vspace{3.5cm}

    \begin{center}
        \textbf{\Large Лабораторная работа №1}\\
        \vspace{0.3cm}
        по курсу «Технологии параллельного программирования»\\
        \vspace{0.8cm}
        \textbf{Освоение программного обеспечения для работы с технологией CUDA.}\\
        \textbf{Примитивные операции над векторами.}
    \end{center}

    \vspace{4cm}

    \begin{flushright}
        Выполнил: Ахметшин Б.Р.\\
        Группа: М8О-403Б-22\\
        Преподаватели: А.Ю. Морозов,\\
        \hspace*{3.6cm} Е.Е. Заяц
    \end{flushright}

    \vfill

    \begin{center}
        Москва, 2026
    \end{center}
\end{titlepage}

\section*{Условие}
Цель работы: Ознакомление и установка программного обеспечения для работы с
программно-аппаратной архитектурой параллельных вычислений (CUDA).

Вариант: Поэлементное умножение векторов.

\section*{Программное и аппаратное обеспечение}
\textbf{GPU}

\begin{tabular}{>{\raggedright\arraybackslash}p{0.50\textwidth}>{\raggedright\arraybackslash}p{0.34\textwidth}}
Имя & NVIDIA-SMI 580.82.07  \\
Compute Capability & 7.5 \\
Объем глобальной памяти & 15637086208 \\
Разделяемая память на блок & 49152 \\
Константная память & 65536 \\
Регистров на блок & 65536 \\
Макс. потоков на блок & (1024, 1024, 64) \\
Макс. размер блока & 2147483647 x 65535 x 65535 \\
Количество мультипроцессоров & 40 \\
\end{tabular}

\newpage
\vspace{0.4cm}
\textbf{CPU}

\begin{tabular}{>{\raggedright\arraybackslash}p{0.50\textwidth}>{\raggedright\arraybackslash}p{0.34\textwidth}}
Архитектура & x86\_64 \\
CPU op-mode(s) & 32-bit, 64-bit \\
Address sizes & 48 bits physical, 48 bits virtual \\
Порядок байт & Little Endian \\
CPU(s) & 12 \\
On-line CPU(s) list & 0-11 \\
ID производителя & AuthenticAMD \\
Имя модели & AMD Ryzen 5 7600 6-Core Processor \\
Семейство ЦПУ & 25 \\
Модель & 97 \\
Потоков на ядро & 2 \\
Ядер на сокет & 6 \\
Сокетов & 1 \\
Степпинг & 2 \\
Frequency boost & enabled \\
CPU max MHz & 5171.0000 \\
CPU min MHz & 545.0000 \\
\end{tabular}

\vspace{0.4cm}
\textbf{RAM}\\
64 GB

\vspace{0.2cm}
\textbf{Data storage}\\
SSD на 1TB

\vspace{0.2cm}
\textbf{ОС}\\
Ubuntu 24.04.4 LTS

\vspace{0.2cm}
\textbf{Среда разработки}\\
VS Code

\vspace{0.2cm}
\textbf{Компилятор}\\
nvcc

\section*{Метод решения}
Для решения задачи используется технология CUDA, позволяющая выполнять
однотипные операции над элементами векторов параллельно на GPU.

Каждый поток CUDA вычисляет результат для одного элемента выходного массива:
\[
c[i] = a[i] \cdot b[i].
\]

Глобальный индекс потока определяется как:
\[
i = \texttt{blockIdx.x * blockDim.x + threadIdx.x}.
\]
Если индекс выходит за пределы диапазона \texttt{[0, n)}, поток не выполняет вычисления.

Перед запуском ядра входные данные копируются с host в device,
после выполнения ядра результат копируется обратно на host.
Для проверки корректности вызовов CUDA API используется макрос \texttt{CSC}.

\section*{Описание программы}
Программа реализована в одном файле \texttt{lab1.cu}.

Основные этапы:
\begin{enumerate}
    \item Считывание размера \texttt{n} и двух входных векторов типа \texttt{double}.
    \item Выделение памяти на GPU: \texttt{cudaMalloc(d\_a, d\_b, d\_c)}.
    \item Копирование входных данных в память устройства: \texttt{cudaMemcpy}.
    \item Запуск CUDA-ядра с параметрами сетки и блока.
    \item Синхронизация устройства и проверка ошибок.
    \item Копирование результата на host и вывод в формате \texttt{\%.10e}.
\end{enumerate}

CUDA-ядро:
\begin{lstlisting}[style=cudastyle]
__global__ void vectorMulKernel(const double *a, const double *b, double *c, int n) {
    int idx = blockIdx.x * blockDim.x + threadIdx.x;
    if (idx < n) {
        c[idx] = a[idx] * b[idx];
    }
}
\end{lstlisting}

Сборка:
\begin{lstlisting}[style=cudastyle]
cd lab1
make
\end{lstlisting}

Запуск:
\begin{lstlisting}[style=cudastyle]
./lab1 < input.txt
\end{lstlisting}

\section*{Результаты}
\textbf{Проверочный пример}

Вход:
\begin{lstlisting}[style=cudastyle]
3
1 2 3
4 5 6
\end{lstlisting}

Выход:
\begin{lstlisting}[style=cudastyle]
4.0000000000e+00 1.0000000000e+01 1.8000000000e+01
\end{lstlisting}

\vspace{0.4cm}
\textbf{Замеры:}

\begin{lstlisting}[basicstyle=\ttfamily\scriptsize,frame=single]
+----+----------+----------------------+----------+----------------+---------------+
| No | n        | CUDA (grid, block)   | CPU, ms  | GPU kernel, ms | GPU total, ms |
+----+----------+----------------------+----------+----------------+---------------+
| 1  | 10000    | <<<    1,   32>>>    | 0.002790 | 0.109888       | 280.195262    |
| 2  | 10000    | <<<  128,  256>>>    |          | 0.093216       | 277.456777    |
| 3  | 10000    | <<< 1024, 1024>>>    |          | 0.128736       | 281.152839    |
| 4  | 100000   | <<<    1,   32>>>    | 0.029381 | 0.065120       | 276.495043    |
| 5  | 100000   | <<<  128,  256>>>    |          | 0.070272       | 284.427908    |
| 6  | 100000   | <<< 1024, 1024>>>    |          | 0.109664       | 281.383196    |
| 7  | 1000000  | <<<    1,   32>>>    | 0.595039 | 0.072160       | 282.053103    |
| 8  | 1000000  | <<<  128,  256>>>    |          | 0.074752       | 285.034988    |
| 9  | 1000000  | <<< 1024, 1024>>>    |          | 0.191328       | 286.654042    |
+----+----------+----------------------+----------+----------------+---------------+
\end{lstlisting}

\section*{Выводы}
\begin{sloppypar}
В лабораторной работе реализована программа поэлементного умножения двух
векторов с использованием технологии CUDA. Алгоритм корректно выполняет
распараллеливание по элементам массива и выводит результат с требуемой точностью.

По результатам замеров видно, что время выполнения самого ядра GPU мало
(порядка $0.065\mbox{--}0.191$ мс), однако полное время GPU в текущей реализации
остается почти постоянным и высоким (около $276\mbox{--}287$ мс) для всех рассмотренных $n$.
Это указывает на доминирование накладных расходов (выделение памяти, копирование
данных, синхронизации и прочие служебные операции), а не времени вычисления ядра.

Сравнение с CPU для заполненных строк показывает, что на данных размера
$10^4$, $10^5$ и $10^6$ текущая GPU-реализация по общему времени уступает CPU,
несмотря на быстрое исполнение самого ядра. Лучшая конфигурация по времени ядра
в данных замерах --- \texttt{<<<1, 32>>>} для $n=10^5$ и $n=10^6$,
а для $n=10^4$ --- \texttt{<<<128, 256>>>}.

Реализация пригодна как базовая учебная версия для освоения CUDA.
Для получения ускорения по полному времени требуется оптимизация обменов
между host и device, а также сокращение накладных расходов запуска.

Почти одинаковое время ядра для конфигураций 1x32 и 128x256 в текущих
замерах объясняется тем, что операция
поэлементного умножения является очень легкой и ограничивается в основном
доступом к памяти и накладными расходами запуска. При этом значения времени
ядра малы и близки к уровню шумов измерения. Дополнительно важно учитывать,
что без grid-stride цикла конфигурация 1x32 не покрывает весь
массив для больших $n$, поэтому такое сравнение не полностью эквивалентно.
\end{sloppypar}

\end{document}
